\documentclass[12pt]{jlreq}
\usepackage{amsmath, amssymb}

\newcommand{\C}{\mathbb{C}}
\newcommand{\R}{\mathbb{R}}
\newcommand{\Z}{\mathbb{Z}}
\newcommand{\Tr}{\operatorname{Tr}}
\newcommand{\Impart}{\operatorname{Im}}
\newcommand{\AF}{\operatorname{AF}}
\newcommand{\EG}{\operatorname{EG}}
\newcommand{\AX}{\operatorname{AX}}
\newcommand{\GL}{\operatorname{GL}}
\newcommand{\Hom}{\operatorname{Hom}}
\newcommand{\Ker}{\operatorname{Ker}}
\newcommand{\Id}{\operatorname{Id}}
\newcommand{\Sym}{\operatorname{Sym}}

\begin{document}

$3$ 次元ユークリッド空間 $\mathbb{R}^3$ に $5$ 個の相異なる平面 $P_1, \ldots, P_5$ が次を満たすように配置されているとする。
\begin{enumerate}
  \item[(i)] 相異なる $2$ 個の平面の共通部分は必ず直線である。
  \item[(ii)] 相異なる $3$ 個の平面の共通部分は必ず $1$ 点からなる。
  \item[(iii)] 相異なる $4$ 個の平面の共通部分は必ず空集合である。
\end{enumerate}
これらの平面の和集合を $X = \bigcup_{i=1}^5 P_i$ とおき，$X$ にユークリッド空間 $\mathbb{R}^3$ の位相の相対位相を与える。このとき，以下の問に答えよ。

\begin{enumerate}
  \item $X$ のオイラー標数を求めよ。
  \item $X$ は単連結であることを示せ。
  \item $X$ の整数係数ホモロジー群を求めよ。
\end{enumerate}

\end{document}